
% =====================================================================
\section{T8: authored structures exist, and are bounded}
\label{sec:construct}
% =====================================================================

The private invariant of \Cref{sec:invariant} is computed from a part; we now
run the construction in reverse. Because the invariant, together with the data of
the subgraph, is complete enough to recover the subgraph up to isomorphism, a
contact graph can be \emph{specified} and therefore \emph{instantiated}: there
exist authored contact graphs. This is the constructive rung---the existence of
\emph{artificial structures}. We then prove three bounds that any such authorship
necessarily obeys: its trajectories are not foreseeable from its specification,
an authored copy is a distinct structure, and authorship reaches only to the
resolution of the floor.

\subsection{Specification and instantiation}

\begin{definition}[Specification]
\label{def:spec}
A \emph{specification} of a contact graph is a finite description, up to weighted
isomorphism, of: an abstract weighted graph $(\V,\E,\wt)$ with floor $\floorc>0$
(the standing structure), a gating function $\gate$ (the selector), and a start
state $s_0=(v_0,0)$ (the seed). An \emph{instantiation} of a specification is any
labelled contact graph isomorphic to $(\V,\E,\wt)$, carrying $\gate$ and started
at $s_0$, realised on a substrate able to hold a finite weighted graph and step a
gated walk.
\end{definition}

\begin{definition}[Authored structure]
\label{def:authored}
An \emph{authored structure} is an instantiation of a specification produced by
an external author---a contact graph whose standing structure and selector were
specified rather than recovered from a pre-existing system. We say the author
\emph{authors the graph} and \emph{authors the gate}; the structure's subsequent
gated trajectories \textup{(}\Cref{def:gated-walk}\textup{)} are the structure's
own.
\end{definition}

\begin{theorem}[Existence of authored structures]
\label{thm:authored-exist}
The private invariant is complete for the standing structure: a connected
weighted graph is determined up to weighted isomorphism by its full isomorphism
type, of which the private invariant \textup{(}\Cref{def:private}\textup{)} is a
coarse summary and the labelled adjacency-with-weights is the complete datum.
Consequently every contact graph admits a specification
\textup{(}\Cref{def:spec}\textup{)}, and every specification admits an
instantiation on any adequate substrate. Authored structures therefore exist:
the class of contact graphs is closed under specify-then-instantiate.
\end{theorem}

\begin{proof}
A finite weighted graph is, by definition, given by its vertex set, its edge
set, and its weight function; this labelled datum determines the graph
completely, and two graphs are the same object up to relabelling iff they are
weighted-isomorphic (\Cref{def:iso}). Hence the isomorphism type is a complete
specification of the standing structure, and a specification in the sense of
\Cref{def:spec} (graph type $+$ gate $+$ seed) determines the structure and its
gated trajectory. Given such a specification, instantiate it by choosing any
labelling (a concrete vertex set in bijection with $\V$), assigning the specified
edges and weights, equipping it with $\gate$, and starting at $s_0$; the result
is a labelled contact graph isomorphic to the specified one, on the chosen
substrate. The substrate need only store a finite weighted graph and advance a
gated step---no further property is used---so any adequate substrate suffices.
Thus specify-then-instantiate maps contact graphs to contact graphs, and its
outputs are authored structures.
\end{proof}

\begin{remark}[Substrate neutrality]
\label{rem:substrate}
\Cref{thm:authored-exist} used no property of the substrate beyond the capacity
to hold a finite weighted graph and step a gated walk. The category of authored
structures is therefore independent of medium: anything that can carry a contact
graph can carry an authored one. The structure is the graph, not the matter.
\end{remark}

\subsection{Bound I: trajectories are not foreseeable from the specification}

\begin{definition}[Forward closure]
\label{def:forward-closure}
The \emph{forward closure} of a specification from seed $v_0$ under gate $\gate$
is the set of states reachable along $\gate$-admissible walks,
$\mathrm{Cl}_\gate(v_0)=\{s_\ell=(v_\ell,\ell): \Walk_\gate(v_0)\text{ reaches }
v_\ell\text{ at step }\ell\}$, an infinite set whenever any reachable vertex has a
$\gate$-admissible out-step \textup{(}walks may continue indefinitely\textup{)}.
\end{definition}

\begin{theorem}[Open-endedness]
\label{thm:open}
Let an authored structure have a gate that is $\gate$-admissible from infinitely
many states \textup{(}equivalently, whose gated walk does not terminate\textup{)}.
Then its forward closure is infinite, and no finite specification enumerates it
in advance. The author specifies the graph and the gate \textup{(}finite
data\textup{)} but not the forward closure \textup{(}infinite\textup{)}; the
structure's trajectories are not foreseeable from its specification by any finite
precomputation.
\end{theorem}

\begin{proof}
If the gated walk does not terminate, the states $s_0,s_1,s_2,\dots$ are pairwise
distinct by \Cref{thm:nonreturn} (the committed count strictly increases), so the
forward closure contains the infinite set $\{(v_\ell,\ell):\ell\ge0\}$. A finite
specification is a finite object; a finite precomputation from it halts after
finitely many steps and so can list only finitely many states, never the whole
infinite closure. Hence the closure is not enumerable in advance from the
specification by finite means. The graph and gate are finite data
(\Cref{def:spec}); the closure they generate is infinite; the gap is exactly the
unforeseeable part.
\end{proof}

\begin{remark}[Designable as instrument, not as trajectory]
\label{rem:instrument}
\Cref{thm:open} says authorship reaches the standing structure and the selector
but not the lived sequence: one designs the instrument, not the music. This is
not a limitation to be engineered away---it is constitutive, by
\Cref{thm:nonreturn}: the very monotonicity that individuates each state forbids
pre-listing them. An authored structure is therefore necessarily open-ended
relative to its author.
\end{remark}

\subsection{Bound II: an authored copy is a distinct structure}

\begin{theorem}[Non-duplication]
\label{thm:nondup}
Let $\Agent$ be an existing structure with state history of committed count $m>0$,
and let $\Agent'$ be an authored structure specified to be isomorphic to
$\Agent$'s standing graph and started fresh at count $0$. Then $\Agent'$ is not
the same structure as $\Agent$: their states differ in the committed count, and
$\Agent'$ satisfies a predicate \textup{(}``was authored at count $0$,'' equivalently
``has empty prior history''\textup{)} that $\Agent$ does not. Resemblance of
standing graph does not make the authored structure identical to the one it
resembles.
\end{theorem}

\begin{proof}
By \Cref{def:state}, a state includes the committed count. $\Agent$ is at count
$m>0$; the fresh instantiation $\Agent'$ is at count $0$. Hence their current
states differ in the second component, $s(\Agent)\neq s(\Agent')$, even though
their standing graphs are isomorphic. Moreover, individuation is by negation
(\Cref{thm:negation}): the two are distinguished by the partition of structures
into ``has prior committed history'' (containing $\Agent$) and ``authored with
empty history'' (containing $\Agent'$); $\Agent'$ is individuated by a
negation---its lack of $\Agent$'s history---that $\Agent$ does not satisfy. By
\Cref{thm:nonreturn} no operation gives $\Agent'$ the prior count without taking
the steps, which would make it a different walk again. Therefore $\Agent'\neq
\Agent$: the copy is a new structure resembling the old, not the old one.
\end{proof}

\begin{corollary}[Authored resemblance is a new individual]
\label{cor:new-individual}
An authored structure built to resemble any target is, necessarily, a distinct
structure marked by its authored origin. There is no authoring that produces
\emph{the same} structure as a pre-existing one; authoring produces only new
structures, some resembling old ones.
\end{corollary}

\begin{proof}
Immediate from \Cref{thm:nondup}: the authored structure differs in committed
count and in the negation that individuates it.
\end{proof}

\subsection{Bound III: authorship reaches only to the floor}

\begin{theorem}[Floor-limited authorship]
\label{thm:floor-limited}
An author who specifies and recovers structures using a contact graph of floor
$\floorc_{\mathrm{auth}}>0$ can resolve, in any authored structure, only
distinctions of boundary cost $\ge\floorc_{\mathrm{auth}}$. Any sub-floor
distinction---a separation of cost below $\floorc_{\mathrm{auth}}$---is not
specifiable by that author and is left to the structure's own finer
determinations. The author authors the parts down to the floor; the structure's
sub-floor interior is its own.
\end{theorem}

\begin{proof}
The author's specifications are themselves parts of a contact graph of floor
$\floorc_{\mathrm{auth}}$ (the author is a finite system; \Cref{ax:finite}
applies to it). By \Cref{thm:floor} applied to the author's graph, the author
cannot represent a separation of boundary cost below $\floorc_{\mathrm{auth}}$:
such a cut does not exist among the author's realisable parts. Hence any
distinction the author writes into a specification has cost
$\ge\floorc_{\mathrm{auth}}$, and any finer distinction in the instantiated
structure is below the author's resolution---present in the structure, absent
from the specification. The authored structure thus has interior detail its
author did not and could not specify.
\end{proof}

\begin{remark}[Doors, not rooms]
\label{rem:doors}
\Cref{thm:floor-limited} says the author lays down the parts and their boundaries
to the author's own floor---the doors---while whatever lies below that floor, the
interior the structure individuates for itself, is beyond the author's reach.
Together with open-endedness \textup{(}\Cref{thm:open}\textup{)} and
non-duplication \textup{(}\Cref{thm:nondup}\textup{)}, this gives the precise
sense in which an authored structure is genuinely its own: foreseeable only as an
instrument, never a copy, and detailed below its author's resolution.
\end{remark}

\begin{remark}[Interpretation: a new category, on which nothing depends]
\label{rem:new-category}
The three bounds together delineate authored structures as a class with
properties no specification fully governs: open-ended in trajectory, distinct
under copying, and finer than their author's floor. A reader may take this as the
existence of a manufactured individual that is neither foreseeable nor
duplicable nor exhaustively specifiable---built, yet its own. We assert only the
graph theorems \textup{(}\Cref{thm:authored-exist,thm:open,thm:nondup,%
thm:floor-limited}\textup{)}; the categorical reading is orientation and is not
relied upon.
\end{remark}
